%! Inverse Trigonometric Functions — Unit 3, Lesson 3.9 — Student Packet Cover Sheet
\documentclass[10pt]{article}
\usepackage{precalculus-article}
\usepackage{precalculus-boxes}
\usepackage{ltablex}
\keepXColumns

\begin{document}

% Full-bleed navy banner
\begin{tikzpicture}[remember picture, overlay]
  \fill[navy] (current page.north west)
    rectangle ([yshift=-0.9in]current page.north east);
\end{tikzpicture}
\vspace{-0.8in}

\begin{center}
  {\color{white}\textbf{\LARGE Precalculus}}\\[4pt]
  {\color{skymid}\large\textbf{Unit 3: Trigonometric and Polar Functions}}\\[6pt]
  {\color{white}\textbf{\large Lesson 3.9 \quad Inverse Trigonometric Functions}}
\end{center}

\vspace{0.22in}
\namedateperiod
\vspace{0.18in}

\begin{learningtargetbox}
\small
\begin{enumerate}[leftmargin=1.5em, itemsep=5pt]
  \item I can explain why the domain of the sine, cosine, and tangent
        functions must be restricted in order to define their inverses.
  \item I can evaluate $\arcsin x$, $\arccos x$, and $\arctan x$ and
        interpret the output as a principal-value angle measure.
  \item I can state the domain and range of each inverse trigonometric
        function and use them to evaluate compositions.
\end{enumerate}
\end{learningtargetbox}

\vspace{0.14in}

\begin{tocbox}
\small
\renewcommand{\arraystretch}{1.5}
\begin{tabularx}{\linewidth}{@{} c l X r @{}}
\rowcolor{sky}
\textbf{\#} & \textbf{Component} & \textbf{Description} & \textbf{Score} \\
\midrule
1 & Warm-Up        & Evaluating trig, domain \& range review         & \blank{1.2cm} \\
2 & Guided Notes   & Restricted domains; evaluating inverse trig      & \blank{1.2cm} \\
3 & Group Activity & Tiers R/A/E --- inverse trig practice            & \blank{1.2cm} \\
4 & Exit Ticket    & Evaluate \& explain domain restriction           & \blank{1.2cm} \\
5 & Homework       & Practice + AP-style MC + composition extension   & \blank{1.2cm} \\
\midrule
  & \multicolumn{2}{r}{\textbf{Total}} & \blank{1.2cm} \\
\end{tabularx}
\end{tocbox}

\end{document}
